%------------------------------------------------------------------------------%
\documentclass[a4paper,12pt,fleqn]{article}

\usepackage{calculo_varias_variaveis-1}

\ShowAnswers

%------------------------------------------------------------------------------%
\begin{document}

\makeHeader{CFVVI}{Prova 1 -- Turma 6}{01/04/2026}

% 01
%------------------------------------------------------------------------------%
\begin{multicols}{2}
\raggedbottom
\exercise{50\%}
  Uma trajetória elíptica é descrita pelas equações paramétricas
  \vspace{-0.5\baselineskip}
  \begin{gather*}
    x(t) = 3 \sin\left(\sqrt{t}\right) - 1  \\
    y(t) = 5 \cos\left(\sqrt{t}\right) + 1  \\
    0 \leq t \leq \pi^2
  \end{gather*}
  \vspace{-1.5\baselineskip}
  \begin{enumerate}[label=\alph*),itemsep=0pt,topsep=0pt]
    \item\!\! [10\%] determine os pontos da curva com
     $t_1 = \frac{\pi^2}{4}$, $t_2 = \pi^2$ e $t_3 = \frac{9\pi^2}{4}$
    \item\!\! [15\%] calcule a expressão para o vetor tangente à curva
    \item\!\! [15\%] avalie o vetor tangente à curva nos pontos $t_1$,  $t_2$ e $t_3$
    \item\!\! [10\%] esboce a curva e os vetores tangentes calculados
  \end{enumerate}
\columnbreak
~\\
\textit{Para esboçar dos vetores, multiplique-os por $\pi$}\\[-2mm]

\begin{questiononly}
  \hfill\includegraphics{axis_xy}
\end{questiononly}
\begin{answeronly}
  \hfill\includegraphics{T6_A_parametrica}
\end{answeronly}
\end{multicols}
\clearpagequestiononly

\begin{answer}
\textbf{a)} \quad
Calculando os pontos
\begin{multicols}{2}
\begin{align*}
  x\left(\frac{\pi^2}{4}\right)
  & = \left(3 \sin\left(\sqrt{t}\right) - 1\right) \evalat{\frac{\pi^2}{4}}{} \\
  & = 3 \sin\left(\sqrt{\frac{\pi^2}{4}}\right) - 1  \\
  & = 3 \sin\left(\frac{\pi}{2}\right) - 1  \\
  & = 3 \times 1 - 1  \\
  & = 2
\end{align*}
\newcolumn
\begin{align*}
  y\left(\frac{\pi^2}{4}\right)
  & = \left(5 \cos\left(\sqrt{t}\right) + 1\right) \evalat{\frac{\pi^2}{4}}{} \\
  & = 5 \cos\left(\sqrt{\frac{\pi^2}{4}}\right) + 1 \\
  & = 5 \cos\left(frac{\pi}{2}\right) + 1 \\
  & = 5 \times 0 + 1 \\
  & = 1
\end{align*}
\end{multicols}

\newpage
\begin{multicols}{2}
\begin{align*}
  x\left(\pi^2\right)
  & = \left(3 \sin\left(\sqrt{t}\right) - 1\right) \evalat{\pi^2}{} \\
  & = 3 \sin\left(\sqrt{\pi^2}\right) - 1  \\
  & = 3 \sin\left(\pi\right) - 1  \\
  & = 3 \times 0 - 1  \\
  & = -1
\end{align*}
\newcolumn
\begin{align*}
  y\left(\pi^2\right)
  & = \left(5 \cos\left(\sqrt{t}\right) + 1\right) \evalat{\pi^2}{} \\
  & = 5 \cos\left(\sqrt{\pi^2}\right) + 1 \\
  & = 5 \cos\left(\pi\right) + 1 \\
  & = 5 (-1) + 1 \\
  & = -4
\end{align*}
\end{multicols}

\begin{multicols}{2}
\begin{align*}
  x\left(\frac{9\pi^2}{4}\right)
  & = \left(3 \sin\left(\sqrt{t}\right) - 1\right) \evalat{\frac{9\pi^2}{4}}{} \\
  & = 3 \sin\left(\sqrt{\frac{9\pi^2}{4}}\right) - 1  \\
  & = 3 \sin\left(\frac{3\pi}{2}\right) - 1  \\
  & = 3 (-1) - 1  \\
  & = -4
\end{align*}
\newcolumn
\begin{align*}
  y\left(\frac{9\pi^2}{4}\right)
  & = \left(5 \cos\left(\sqrt{t}\right) + 1\right) \evalat{\frac{9\pi^2}{4}}{} \\
  & = 5 \cos\left(\sqrt{\frac{9\pi^2}{4}}\right) + 1 \\
  & = 5 \cos\left(\frac{3\pi}{2}\right) + 1 \\
  & = 5 \times 0 + 1 \\
  & = 1
\end{align*}
\end{multicols}


\textbf{b)} \quad
O vetor tangente é
\(
  v(t) = \Vetor{\frac{dx}{dt}}{\frac{dy}{dt}}
\)

\begin{multicols}{2}
\begin{align*}
  \frac{dx}{dt}
  & = \frac{d}{dt}\left[3 \sin\left(\sqrt{t}\right) - 1\right] \\[2mm]
  & = 3 \frac{d}{dt}\left[\sin\left(\sqrt{t}\right)\right] \\[2mm]
  & = 3 \cos\left(\sqrt{t}\right)\;\frac{d}{dt}\left[\sqrt{t}\right] \\[2mm]
  & = 3 \cos\left(\sqrt{t}\right)\frac{1}{2\sqrt{t}} \\[2mm]
  & = \frac{3 \cos\left(\sqrt{t}\right)}{2\sqrt{t}}
\end{align*}
\newcolumn
\begin{align*}
  \frac{dy}{dt}
  & = \frac{d}{dt}\left[5 \cos\left(\sqrt{t}\right) + 1\right]   \\[2mm]
  & = 5\frac{d}{dt}\left[\cos\left(\sqrt{t}\right)\right]   \\[2mm]
  & = -5\sin\left(\sqrt{t}\right)\frac{d}{dt}\left[\sqrt{t}\right]   \\[2mm]
  & = -5\sin\left(\sqrt{t}\right)\frac{1}{2\sqrt{t}} \\[2mm]
  & = \frac{-5\sin\left(\sqrt{t}\right)}{2\sqrt{t}}
\end{align*}
\end{multicols}
\(
  v(t)
  = \Vetor%
    {\frac{3 \cos\left(\sqrt{t}\right)}{2\sqrt{t}}}
    {\frac{-5\sin\left(\sqrt{t}\right)}{2\sqrt{t}}}
  = \frac{1}{2\sqrt{t}}\Vetor%
    { 3\cos\left(\sqrt{t}\right)}
    {-5\sin\left(\sqrt{t}\right)}
\)

\bigskip
\textbf{c)} \quad
Avaliando o vetor tangente nos pontos
$t_1 = \frac{\pi^2}{4}$, $t_2 = \pi^2$ e $t_3 = \frac{9\pi^2}{4}$
\begin{align*}
  v\left(\frac{\pi^2}{4}\right)
  & = \frac{1}{2\sqrt{t}}\Vetor%
      { 3\cos\left(\sqrt{t}\right)}
      {-5\sin\left(\sqrt{t}\right)}
      \evalat{\frac{\pi^2}{4}}{}
  = \frac{1}{2\sqrt{\frac{\pi^2}{4}}}\Vetor%
      { 3\cos\left(\sqrt{\frac{\pi^2}{4}}\right)}
      {-5\sin\left(\sqrt{\frac{\pi^2}{4}}\right)}
  \\
  & = \frac{1}{2\frac{\pi}{2}}\Vetor%
      { 3\cos\left(\frac{\pi}{2}\right)}
      {-5\sin\left(\frac{\pi}{2}\right)}
  = \frac{1}{\pi}\Vetor%
      { 3\cos\left(\frac{\pi}{2}\right)}
      {-5\sin\left(\frac{\pi}{2}\right)}
  = \frac{1}{\pi}\vetor%
      { 3\times 0}
      {-5\times 1}
  = \frac{1}{\pi}\vetor%
      { 0}
      {-5}
\end{align*}
\begin{align*}
  v\left(\pi^2\right)
  & = \frac{1}{2\sqrt{t}}\Vetor%
      { 3\cos\left(\sqrt{t}\right)}
      {-5\sin\left(\sqrt{t}\right)}
      \evalat{\pi^2}{}
  = \frac{1}{2\sqrt{\pi^2}}\Vetor%
      { 3\cos\left(\sqrt{\pi^2}\right)}
      {-5\sin\left(\sqrt{\pi^2}\right)}
  \\
  & = \frac{1}{2\pi}\Vetor%
      { 3\cos\left(\pi\right)}
      {-5\sin\left(\pi\right)}
  = \frac{1}{2\pi}\vetor%
      { 3(-1)}
      {-5\times 0}
  = \frac{1}{2\pi}\vetor%
      {-3}
      {0}
\end{align*}
\begin{align*}
  v\left(\frac{9\pi^2}{4}\right)
  & = \frac{1}{2\sqrt{t}}\Vetor%
      { 3\cos\left(\sqrt{t}\right)}
      {-5\sin\left(\sqrt{t}\right)}
      \evalat{\frac{9\pi^2}{4}}{}
  = \frac{1}{2\sqrt{\frac{9\pi^2}{4}}}\Vetor%
      { 3\cos\left(\sqrt{\frac{9\pi^2}{4}}\right)}
      {-5\sin\left(\sqrt{\frac{9\pi^2}{4}}\right)}
  \\
  & = \frac{1}{2\frac{3\pi}{2}}\Vetor%
      { 3\cos\left(\frac{3\pi}{2}\right)}
      {-5\sin\left(\frac{3\pi}{2}\right)}
  = \frac{1}{3\pi}\Vetor%
      { 3\cos\left(\frac{3\pi}{2}\right)}
      {-5\sin\left(\frac{3\pi}{2}\right)}
  = \frac{1}{3\pi}\vetor%
      { 3\times 0}
      {-5(-1)}
  = \frac{1}{3\pi}\vetor%
      {0}
      {5}
\end{align*}
\end{answer}

% 02
%------------------------------------------------------------------------------%
\exercise{50\%}
Considerando a função
\(
  f(x, y) = \sqrt{\ln\left(10-x^2+y^2\right)}
\)
\begin{enumerate}[label=(\alph*),itemsep=0pt]
  \item \ [15\%] determine o domínio de $f$
  \item \ [15\%] esboce o domínio
  \item \ [10\%] determine a curva de nível que passa pelo ponto \((0, 0)\)
  \item \ [10\%] esboce a curva de nível que passa pelo ponto \((0, 0)\)
\end{enumerate}
\begin{questiononly}
  \begin{multicols}{2}
    \;\; Domínio
    \vspace{-0.3\baselineskip}
    \begin{center}
      \includegraphics{axis_xy}
    \end{center}
    \;\; Curva de nível
    \vspace{-0.3\baselineskip}
    \begin{center}
      \includegraphics{axis_xy}
    \end{center}
  \end{multicols}
\end{questiononly}
\begin{answeronly}
  \begin{multicols}{2}
    \;\; Domínio
    \vspace{-0.3\baselineskip}
    \begin{center}
      \includegraphics{T6_A_dominio}
    \end{center}
    \;\; Curva de nível
    \vspace{-0.3\baselineskip}
    \begin{center}
      \includegraphics{T6_A_curva_nivel}
    \end{center}
  \end{multicols}
\end{answeronly}

\begin{answer}
\begin{multicols}{2}
  \textbf{a)} \quad
  Para calcular a raiz precisamos que
  \begin{align*}
    \ln\left(10-x^2+y^2\right)        & \geq 0 \\
    10-x^2+y^2                        & \geq 1 \\
    -x^2+y^2                          & \geq 9 \\
    x^2 - y^2                         & \leq 9 \\
    \frac{x^2}{3^2} - \frac{y^2}{3^2} & \leq 1
  \end{align*}
  ou seja, os pontos ``entre os lados'' da hipérbole.
  Assim, o domínio é
  \[
    D_f = \{(x, y) \in \R^2 \st x^2 - y^2 < 9\}
  \]

  \bigskip
  \textbf{c)} \quad
  A curva de nível que passa pelo ponto $(0, 0)$ tem o valor
  \begin{align*}
    c
    & = f(0, 0) \\
    & = \sqrt{\ln\left(10 -x^2 +y^2 \right)} \evalat{(0, 0)}{} \\
    & = \sqrt{\ln\left(10 -0^2 +0^2 \right)} \\
    & = \sqrt{\ln 10}
  \end{align*}
  A curva de nível associada é
  \begin{align*}
    f(x, y)                           & = c             \\[2mm]
    \sqrt{\ln\left(10-x^2+y^2\right)} & = \sqrt{\ln 10} \\[2mm]
    \ln\left(10-x^2+y^2\right)        & = \ln 10        \\[2mm]
    10 - x^2 + y^2                    & = 10            \\[2mm]
    - x^2 + y^2                       & = 0             \\[2mm]
    y^2                               & = x^2           \\[2mm]
    y                                 & = \pm x
  \end{align*}
  A curva de nível corresponde as retas
  \[
    y = x
    \qquad\text{e}\qquad
    y = -x
  \]
\end{multicols}
\end{answer}

%------------------------------------------------------------------------------%
\end{document}
%------------------------------------------------------------------------------%
